\documentclass[10pt]{article}
\usepackage[T1]{fontenc}
\usepackage[utf8]{inputenc}
\usepackage[french]{babel}
\usepackage{amsmath,amssymb,mathrsfs,mathpazo}
\usepackage{geometry}
\geometry{a4paper,margin=2cm}
\usepackage{graphicx,array,tabularx,tabularray,colortbl,multirow,multicol,booktabs,enumitem,pifont,ulem,fancybox,diagbox,cancel}
\usepackage{pstricks,pstricks-add,pst-plot,pst-tree,pst-node,pst-text,pst-func,pst-math,pst-bspline,pst-3dplot,pst-eucl,pst-solides3d}
\let\pscalculate\undefined
\usepackage{pst-bezier}
\usepackage{tikz,pgfplots,tkz-tab}
\usetikzlibrary{arrows,arrows.meta,patterns,calc,positioning,trees,decorations.markings,shapes.geometric,shapes.misc,shapes.arrows,intersections,decorations.pathmorphing,decorations.pathreplacing,backgrounds,shadows,babel}
\pgfplotsset{compat=1.18}
\usepgfplotslibrary{fillbetween}
\usepackage[np]{numprint}
\usepackage{listings,makecell,pas-tableur,twoopt}
\let\barre\undefined
\usepackage[pstricks]{tablvar}
\definecolor{darkgreen}{RGB}{0,100,0}
\providecommand{\red}{\color{red}}
\providecommand{\blue}{\color{blue}}
\providecommand{\green}{\color{green}}
\providecommand{\black}{\color{black}}
\providecommand{\white}{\color{white}}
\providecommand{\gray}{\color{gray}}
\providecommand{\cyan}{\color{cyan}}
\providecommand{\magenta}{\color{magenta}}
\providecommand{\yellow}{\color{yellow}}
\providecommand{\violet}{\color{violet}}
\providecommand{\orange}{\color{orange}}
\providecommand{\euro}{\texteuro}
\providecommand{\EUR}{\texteuro}
\providecommand{\ieme}{\textsuperscript{e}}
\providecommand{\degres}{\ensuremath{^\circ}}
\providecommand{\degree}{\ensuremath{^\circ}}
\providecommand{\vv}[1]{\overrightarrow{#1}}
\providecommand{\Coord}[1]{#1}
\providecommand{\decosix}{}
\providecommand{\starredbullet}{\ensuremath{\bullet}}
\pagestyle{empty}
\setlength{\parindent}{0pt}
\usepackage{adjustbox}\DeclareUnicodeCharacter{2236}{:}\DeclareUnicodeCharacter{2190}{\ensuremath{\leftarrow}}
\providecommand{\R}{\mathbb{R}}
\providecommand{\N}{\mathbb{N}}
\providecommand{\Z}{\mathbb{Z}}
\providecommand{\Q}{\mathbb{Q}}
\providecommand{\D}{\mathbb{D}}
\providecommand{\C}{\mathbb{C}}
\providecommand{\e}{{\mathrm{e}}}
\providecommand{\vect}[1]{\overrightarrow{#1}}
\providecommand{\vectt}[1]{\overrightarrow{\mathrm{#1}}}
\providecommand{\vvt}[1]{\overrightarrow{\mathrm{#1}}}
\providecommand{\np}[2][]{{#2~#1}}
\providecommand{\vv}[1]{{\overrightarrow{#1}}}
\renewcommand{\d}{{\mathrm{d}}}
\providecommand{\degres}{{^\circ}}
\providecommand{\degre}{{^\circ}}
\providecommand{\blue}{\color{blue}}
\providecommand{\red}{\color{red}}
\providecommand{\enskip}{\quad}
\providecommand{\llbracket}{\lbrack\!\lbrack}
\providecommand{\rrbracket}{\rbrack\!\rbrack}
\providecommand{\euro}{{\text{€}}}
\providecommand{\no}{{\text{n°}}}
\providecommand{\interval}[2]{{[#1;#2]}}
\providecommand{\intervFO}[2]{{[#1;#2[}}
\providecommand{\dotfill}{\ldots}

\providecommand{\ds}{\displaystyle}
\providecommand{\textding}[1]{#1}
\providecommand{\Oij}{\(\mathrm{O};\vec{\imath},\vec{\jmath}\)}
\providecommand{\Oijk}{\(\mathrm{O};\vec{\imath},\vec{\jmath},\vec{k}\)}
\providecommand{\Ouv}{\(\mathrm{O};\vec{u},\vec{v}\)}
\providecommand{\intervFF}[2]{\ensuremath{[#1;#2]}}
\renewcommand{\euro}{\ifmmode\text{\texteuro}\else\texteuro\fi}
\DeclareUnicodeCharacter{00B0}{\ifmmode^{\circ}\else\textdegree\fi}
\DeclareUnicodeCharacter{20AC}{\ifmmode\text{\texteuro}\else\texteuro\fi}
% Énoncés : Éducation nationale. Transcription et corrigés : APMEP, Denis Vergès et contributeurs. Découpage et classement : CoopMaths, Sébastien Lozano, Éric Elter et contributeurs. Adaptation de la présentation : Maths Explorer.
% CC BY-SA : https://coopmaths.fr/www/about/
% https://coopmaths.fr/alea/static/bac/2025/tex/bac_2025_11_sujet2_ameriquesud_3.tex
\begin{document}
\noindent\textit{APMEP / CoopMaths. CC BY-SA. Présentation adaptée par Maths Explorer.}\par\medskip
\medskip


On considère les suites $(v_n)$ et $(w_n)$ définies pour tout entier naturel $n$
par :
\[ \left\{\begin{array}{l c l}
v_0 &=& \ln (4)\\
v_{n+1} &=& \ln \left(- 1 +2\e^{v_n}\right)
\end{array}\right. \qquad \text{et} \qquad w_n = - 1 + \e^{v_n}.\]

On admet que la suite $(v_n)$ est bien définie et strictement positive.

\medskip

\begin{minipage}{0.6\linewidth}
\begin{enumerate}
\item Donner les valeurs exactes de $v_1$ et $w_0$.
\item 
	\begin{enumerate}
		\item Une partie d'une feuille de calcul où figurent les indices et
les termes des suites $(v_n)$ et $(w_n)$ est reproduite ci-contre.
		
Parmi les trois formules ci-dessous, choisir la formule qui,
saisie dans la cellule B3 puis recopiée vers le bas, permettra
d'obtenir les valeurs de la suite $(v_n)$ dans la colonne B.

\begin{tabular}{|l |l|}\hline
Formule 1 &LN($-1 + 2$ * EXP(B2))\\ \hline
Formule 2 &= LN($-1 + 2$ * EXP(B2))\\ \hline
Formule 3 &= LN($-1 + 2$ * EXP(A2))\\ \hline
\end{tabular}
		\item Conjecturer le sens de variation de la suite $(v_n)$.
		\item À l'aide d'un raisonnement par récurrence, valider votre
conjecture concernant le sens de variation de la suite $(v_n)$.
	\end{enumerate}
\end{enumerate}
\end{minipage} \hfill 
\begin{minipage}{0.36\linewidth}
{\small\begin{tabular}{c |c|c|c}\hline
	&A &B &C\\ \hline
1 	&$n$&$v_n$& $w_n$\\ \hline
2&0 & 1,38629436 &3\\ \hline
3& 1& 1,94591015 &6\\ \hline
4& 2& 2,56494936 &12\\ \hline
5& 3& 3,21887582 &24\\ \hline
6& 4& 3,8918203 &48\\ \hline
7& 5& 4,57471098 &96\\ \hline
8& 6& 5,26269019 &192\\ \hline
9& 7& 5,95324333 &384\\ \hline
10& 8& 6,6450909 7 &768\\ \hline
11& 9& 7,33758774 &1536\\ \hline
12& 10& 8,03040956 &3072\\ \hline
13& 11& 8,72339402 &6144\\ \hline
14& 12& 9,41645983 &12288\\ \hline
15& 13& 10,1095663 &24576\\ \hline
16& 14& 10,8026932 &49152\\ \hline
17& 15& 11,4958302&98304\\ \hline
18& 16& 12,1889723&196608\\ \hline
19& 17&12,8821169& 393216\\ \hline
\end{tabular}}
\end{minipage}

\begin{enumerate}[start=3]
\item
	\begin{enumerate}
		\item Démontrer que la suite $(w_n)$ est géométrique.
		\item En déduire que pour tout entier naturel $n, \: v_n = \ln \left(1 + 3 \times 2^n\right)$.
		\item Déterminer la limite de la suite $(v_n)$.
	\end{enumerate}
\item Justifier que l'algorithme suivant écrit en langage Python renvoie un résultat quel que soit le choix de la valeur du nombre S.

\begin{center}
\begin{ttfamily}
\begin{tabular}{|l|}\hline
from math import*\\
def seuil(S):\\
\qquad V=ln(4)\\
\qquad n=0\\
\qquad while V < S :\\
\qquad  \quad n=n+1\\
\qquad  \quad V=ln(2*exp(V)-1)\\
\qquad return(n)\\ \hline
\end{tabular}
\end{ttfamily}
\end{center}
\end{enumerate}

\end{document}
